\section{Feature Analysis and Importance}
\label{sec:importance}

Two analyses answer the main research question. SHAP attributes each
prediction to individual features, and a controlled ablation measures what the
school-level data adds over municipal socioeconomics alone.

\textbf{What the models actually rely on.} Figure~\ref{fig:shap} ranks the nine
features confirmed for all four targets, plus municipal \texttt{cluster} for
comparison, by mean absolute SHAP value as a share of each target's total so the
outcomes sit on one scale. These values come from a RandomForest refitted on all
rows using the same Boruta-selected features per target, matching the family
Table~\ref{tab:leaderboard} selects, so they explain that model rather than
held-out predictions, and they report magnitude, not direction.

The \texttt{nurture\_quintile} leads every target, taking 19.7\% to 36.2\% of
total attribution. This is the Ministry's nurture (\emph{tipuach}) index recorded
\emph{per school} on a 1--5 scale where \textbf{higher means more disadvantaged}:
mean combined grade falls monotonically from 85.3 at quintile 1 to 74.1 at
quintile 5. It is missing for 11.9\% of school-years, which leave the
complete-case sample. Municipal \texttt{cluster} never exceeds 5.3\% and is
dropped entirely for English average grade. Both measure socioeconomic standing,
so this is not a finding that socioeconomics do not matter but one about
\emph{granularity}: schools inside one municipality differ enough that the
town-level average conceals most of the effect.

Beyond that index the leading features are institutional. The transport budget
line ranks second for every target, but needs care: it is negative for 86\% of
institutions in the raw Ministry file and 83\% in the complete-case modelling
sample, indicating a net accounting adjustment rather than gross spending, so
its meaning needs confirmation against the Ministry's data dictionary. Class
size and school size follow, and school size is especially strong for advanced
Math participation (16.5\%), consistent with a large cohort being what makes
opening a 5-unit class viable.

\begin{figure}[H]
\centering
\includegraphics[width=0.92\linewidth]{section7/shap_core_ranking.png}
\caption{Mean $|$SHAP$|$ importance as a share of each target's total, from a
RandomForest refit matching the selected champion. The school-level nurture
index leads all four outcomes (19.7\% to 36.2\%), while municipal cluster never
exceeds 5.3\% and is not selected at all for English average grade.}
\label{fig:shap}
\end{figure}

\textbf{How much the school-level data adds.} SHAP ranks features inside a model
but cannot say what the extra data was worth. So each target is modelled twice on
\emph{identical rows} under the same nested protocol: once on municipal features
alone (cluster, population, locality type, year) and once on the full candidate
set. Only the available information differs, so the gap measures incremental
predictive value (Table~\ref{tab:ablation}).
\phantomsection\label{sec:secondary}

\begin{table}[H]
\centering
\footnotesize
\renewcommand{\arraystretch}{1.3}
\caption{Ablation study under the same nested protocol: outer-fold $R^2$
(mean $\pm$ sd) from municipal features alone, and the gain from adding the
school-level set on identical rows. The full-set scores are the Random Forest
column of Table~\ref{tab:leaderboard} and so are not repeated here.}
\label{tab:ablation}
\begin{tabular}{@{}>{\raggedright\arraybackslash}p{3.2cm}rccc@{}}
\toprule
\textbf{Target} & \textbf{Rows} & \textbf{Municipal only} & \textbf{$\Delta R^2$ from school data} & \textbf{Folds improved} \\
\midrule
Math grade              & 2,995 & $0.109\!\pm\!0.019$ & $\mathbf{+0.311\!\pm\!0.040}$ & 5/5 \\
English grade           & 2,963 & $0.196\!\pm\!0.028$ & $\mathbf{+0.256\!\pm\!0.034}$ & 5/5 \\
Math 5-unit             & 3,221 & $0.055\!\pm\!0.033$ & $\mathbf{+0.340\!\pm\!0.039}$ & 5/5 \\
English 5-unit          & 3,226 & $0.229\!\pm\!0.050$ & $\mathbf{+0.328\!\pm\!0.043}$ & 5/5 \\
\midrule
\textbf{Mean} & & \textbf{0.147} & $\mathbf{+0.309}$ & \textbf{20/20} \\
\bottomrule
\end{tabular}
\end{table}

Municipal features alone predict little, with advanced Math participation close
to unpredictable from the town. Adding the school-level set raises every target
to between 0.395 and 0.557, a mean gain of $+0.309$ holding in \textbf{every
outer fold}, so it is not an artefact of a favourable split. The result is
associational: incremental predictive value, not a causal effect of school
resources.

\textbf{Subject resilience and low-status overachievers.} The secondary questions
use the municipal gradient alone, holding institutional features out so they read
as a clean measure of socioeconomic disparity (Table~\ref{tab:secondary}).
Comparing cluster~2 against cluster~9, Mathematics is the more resilient subject:
a slightly smaller achievement gap and a participation gap roughly a third of
English's. Overachievers are clusters~1--4 schools whose average grade reaches
the \emph{median} of clusters~8--10 (84.90). Of 460 such schools, \textbf{87
(18.9\%)} clear that bar, and on Welch two-sample tests their advanced-track
participation is roughly 2.7$\times$ their peers' in Math and 2.6$\times$ in
English.

\begin{table}[H]
\centering
\footnotesize
\renewcommand{\arraystretch}{1.25}
\caption{Secondary questions. Resilience contrasts cluster~2 with cluster~9.
Overachievers are clusters~1--4 schools at or above the clusters~8--10 median
grade (84.90). The comparison group is the remaining low-cluster schools.}
\label{tab:secondary}
\begin{tabular}{@{}>{\raggedright\arraybackslash}p{5.0cm}rrr@{}}
\toprule
\textbf{(A) Resilience, cluster 2 $\rightarrow$ 9} & \textbf{Math} & \textbf{English} & \\
Grade gap (points), Cohen's $d$ & 6.18 (0.91) & 6.45 (1.16) & \\
5-unit participation gap & \textbf{0.115} & 0.367 & \\
\midrule
\textbf{(B) Overachievers, 87 of 460} & \textbf{Over} & \textbf{Peers} & \textbf{$p$} \\
Math 5-unit participation & 0.132 & 0.049 & $1.4\!\times\!10^{-5}$ \\
English 5-unit participation & 0.468 & 0.181 & $1.4\!\times\!10^{-11}$ \\
\bottomrule
\end{tabular}
\end{table}
